\section{二、问题分析}\label{ux4e8cux95eeux9898ux5206ux6790}

\subsection{2.1
数据探索与特征剖析}\label{ux6570ux636eux63a2ux7d22ux4e0eux7279ux5f81ux5256ux6790}

附件提供了2019年1月2日至2022年5月21日期间共22,453条订单记录，涵盖284种物料，具有以下显著特征：
1.
\textbf{时间跨度与周度离散化}：以2019年1月2日（周三）至1月6日（周日）为第1周，其后每周自周一至周日独立成周，全样本精准跨越177周。
2. \textbf{多品种小批量的间歇性与块状需求（Intermittent \& Lumpy
Demand）}： 根据经典 Syntetos-Boylan 分类准则，定义平均需求间隔
\(\text{ADI} = \frac{1}{N-1}\sum (t_k - t_{k-1})\) 与需求平方变异系数
\(CV^2 = (\sigma_D/\mu_D)^2\)。经数据画像计算，全样本284种物料中，92.6\%的物料属于高方差与高离散间隔的``块状''（Lumpy）或``间歇''（Intermittent）需求，传统的标准高斯分布时序模型易失效，必须引入
Croston、SBA（Syntetos-Boylan Approximation）等专业间歇性时序算法。 3.
\textbf{物料价值异质性极其剧烈}： 物料销售单价从最低 38.16 元到最高
18,676.34
元，跨越近500倍；且部分单价极高物料（如6004020900，单价6331.49元）多余存货1件将占用巨额现金流，而低单价物料缺货1件将造成客户满意度骤降。

\subsection{2.2
四大子问题的逻辑递进关系}\label{ux56dbux5927ux5b50ux95eeux9898ux7684ux903bux8f91ux9012ux8fdbux5173ux7cfb}

各子问题呈现由静态画像到动态排产、由单目标履约到多目标资金权衡、由单步滞后到任意多步延迟的严密递进链条：

\begin{enumerate}
\def\labelenumi{\arabic{enumi}.}
\tightlist
\item
  \textbf{问题 1
  分析}：解决``管什么与怎么看未来''。构建多维评价指标矩阵，采用客观赋权的熵权TOPSIS模型筛选最具代表性的6种物料；在1\textasciitilde100周历史期内以滚动交叉验证遴选最优预测器。
\item
  \textbf{问题 2 分析}：解决``单步滞后下的动态安全库存排产''。生产提前期
  \(L=1\)
  意味着本周决策下周到货。排产模型需通过跟踪预测残差波动，自适应设定动态安全库存
  \(SS = z \cdot \sigma_e\)，在闭环反馈中精准抑制缺货率，达成服务水平
  \(\ge 85\%\)。
\item
  \textbf{问题 3 分析}：解决``单价异质性下的资金-履约双目标 Pareto
  优化''。将单价 \(c_i\) 转化为资金持有成本率 \(h\) 与缺货惩罚
  \(\gamma\)，构建多目标权衡前沿，对高单价物料实施精益防呆、低单价物料实施充足缓冲。
\item
  \textbf{问题 4 分析}：解决``任意时滞 \(k\ge 2\)
  管道在途库存的控制论推广''。当 \(L=k\ge 2\)
  时，系统中存在多期处于运输或加工阶段的``在途库存''（Pipeline
  Inventory）。若决策层仅根据当前在库库存排产，必将导致超量订货与震荡。需构建涵盖在途库存补偿的离散状态空间控制律。
\end{enumerate}
